\documentclass[addpoints]{exam}
% \documentclass[addpoints,12pt]{exam}

\usepackage[slovene]{babel}
\usepackage{amsfonts}
\usepackage{amssymb}
\usepackage{amsmath}
\usepackage{float}
\usepackage{graphicx}
\usepackage{enumitem}
\usepackage{color}
\usepackage[gen]{eurosym}
\usepackage{newunicodechar}
\newunicodechar{€}{\euro}


%Komentar naslednje vrstice glede na verzijo testa -- rešitve ali prostor za odgovore:
% \printanswers

% \linespread{2}


\pointsinrightmargin
\marginbonuspointname{}


\renewcommand{\solutiontitle}{\noindent\textbf{Rešitve in točkovanje:}\par\noindent}

\colorgrids
\definecolor{GridColor}{gray}{0.7}
\setlength{\gridsize}{7mm}

\extrawidth{5mm}
\setlength{\rightpointsmargin}{1.5cm}

\begin{document}

\runningfooter{}{}{}

\begin{center}
    \fbox{\fbox{\parbox{15cm}{\centering
    Peto pisno ocenjevanje znanja -- ponovno \\ 1. letnik -- test F \\ 4. junij 2025 \\ (Koordinatni sistem v ravnini, Funkcija, Premica)}}}
\end{center}

\vspace{0.1in}
\makebox[\textwidth]{Ime in priimek, oddelek:\enspace\hrulefill}


\begin{center}
    \chqword{Naloga}
    \chpword{Možne točke}
    \chbpword{Dodatne točke}
    \chsword{Dosežene točke}
    \chtword{Skupaj}  
    % \combinedpointtable[h][questions]
    \combinedgradetable[h][questions]

\end{center}

\vspace{0.1in}


\begin{questions}

    % 1. vprašanje

    \question
    Narišite dano množico točk v ravnini.
    \begin{parts}
        \part[4]
        $[-2,\infty)\times [-1,4)$

            \begin{solution}[7cm]
                ????
            \end{solution}

        \part[4]
        $\left\{(x,y)\in\mathbb{R}^2 \mid (y\leq 3)\land (y=3x-2)\right\} $

            \begin{solution}[5.5cm]
                ????
            \end{solution}

    \end{parts}
    
    
    \newpage

    % 2. vprašanje
    \question[4]
    Izračunajte obseg trikotnika $\triangle ABC$ z oglišči v točkah $A(1,-4)$, $B(4,5)$ in $C(-3,4)$.
    
    
    \begin{solution}[5cm]
        ???
    \end{solution}

        


    % 3. vprašanje
    \question
    Dana je premica z enačbo $y=-2x+3$, ki seka koordinatni osi v točkah $M$ in $N$.

    \begin{parts}
        \part[4]
        Narišite premico.

            \begin{solution}[6cm]
                ???
            \end{solution}

        \part[4]
        Določite segmentno obliko enačbe pravokotnice k dani premici, ki poteka skozi točko $(-2,1)$. % Premico tudi narišite.

            \begin{solution}[5cm]
                ???
            \end{solution}

        \part[4]
        Naj bo $S$ razpolovišče daljice $MN$. Izračunajte razdaljo med točko $S$ in koordinatnim izhodiščem.

            \begin{solution}[2cm]
                ???
            \end{solution}

    \end{parts}
    
    \newpage


    % 5. vprašanje
    \question
    Dana je funkcija $$f(x)=\begin{cases}
    -2x-2, & x\leq 0 \\ -3, & 1<x<2 \\ -3x+3, & x\geq 2
    \end{cases}.$$ 
    
    \begin{parts}
        \part[5]
        Narišite dano funkcijo.

            \begin{solution}[9cm]
                ???
            \end{solution}

        \part[3]
        Določite zalogo vrednosti in izračunajte ničli funkcije.

            \begin{solution}[7cm]
                ???
            \end{solution}

        % \part[2]
        % Določite in utemeljite, ali je funkcija surjektivna. 
        
        %     \begin{solution}[3cm]
        %         ???
            % \end{solution}

        \part[1]
        Določite območje, kjer je funkcija strogo padajoča.
        
            \begin{solution}[2cm]
                ???
            \end{solution}

    \end{parts}
    
    


    \newpage 


    % 4. vprašanje
    \question
        Določite vrednost parametra $a$, da bo premica $ax+2y+a+5=0$:
        \begin{parts}
            \part[2] 
            potekala skozi točko $(-8,1)$,

                \begin{solution}[6cm]
                    ???
                \end{solution}


            % \part[3]
            % padajoča,
            
            %     \begin{solution}[5cm]
            %         ???
            %     \end{solution}

            \part[3] 
            vzporedna simetrali lihih kvadrantov,
                        
                \begin{solution}[6cm]
                    ???
                \end{solution}

                
            \bonuspart[4]
            \textit{Dodatna naloga:} sekala premico $y=2x+1$ pri $x=-2$.

        \end{parts}
    



\end{questions}



\end{document}